\chapter{Estimação de Parâmetros}
\label{cap:05}

Modelos matemáticos de sistemas biológicos contêm parâmetros --- taxas de reação,
constantes de afinidade, coeficientes de degradação, capacidades de transporte --- cujos
valores devem ser determinados a partir de dados experimentais. Esse processo é chamado de
\emph{estimação de parâmetros} e constitui um dos problemas centrais da biologia de
sistemas quantitativa. Sem parâmetros adequados, mesmo o modelo estruturalmente correto
produzirá previsões sem valor.

O desafio fundamental é que estimar parâmetros é o inverso de simular: dado um modelo e
parâmetros, simulação produz dados; dada uma trajetória de dados e um modelo, estimação
produz parâmetros. O problema inverso é matematicamente mais difícil que o direto. É
frequentemente \emph{mal-posto} no sentido de Hadamard: pode não ter solução única, pode ser
sensível a perturbações nos dados, e pode ter múltiplos conjuntos de parâmetros que ajustam
igualmente bem os dados observados. Compreender essas dificuldades é tão importante quanto
conhecer os algoritmos de otimização.

Este capítulo cobre os três pilares da estimação de parâmetros: a \emph{formulação
matemática} do problema (como definir o que é ``ajuste adequado''), os \emph{algoritmos de
otimização} (como encontrar os parâmetros que minimizam o desajuste), e a
\emph{análise estatística} do ajuste (como quantificar incerteza, identificabilidade e
capacidade preditiva do modelo).

\section{O Problema Inverso}

\begin{definicaoenv}[Estimação de Parâmetros]
Dado um modelo matemático $\hat{y}(t; \boldsymbol{\theta})$ e um conjunto de observações
experimentais $\{(t_i, y_i)\}_{i=1}^n$, o problema de estimação de parâmetros consiste em
encontrar o vetor $\boldsymbol{\theta}^*$ que minimiza uma medida de discrepância entre as
previsões do modelo e os dados:
\begin{equation}
\boldsymbol{\theta}^* = \arg\min_{\boldsymbol{\theta}} J(\boldsymbol{\theta})
\end{equation}
onde $J(\boldsymbol{\theta})$ é a \emph{função objetivo} (ou função de custo).
\end{definicaoenv}

A diferença entre o problema direto e o inverso é pedagogicamente iluminadora. No problema
direto, conhecemos os parâmetros --- por exemplo, $k_1 = 0.5$ e $k_2 = 0.1$ para o modelo
$\dot{x} = k_1 - k_2 x$ --- e calculamos a trajetória $x(t)$ por integração numérica. O
resultado é único e determinístico. No problema inverso, observamos $x(t)$ em pontos
discretos e com ruído experimental, e devemos inferir $k_1$ e $k_2$. Esse processo é não
único (múltiplos pares $(k_1, k_2)$ podem gerar trajetórias igualmente próximas dos dados),
sensível ao ruído (pequenas perturbações nos dados podem causar grandes variações nos
parâmetros estimados), e computacionalmente custoso (requer múltiplas simulações do modelo).

Os principais desafios práticos na estimação de parâmetros em modelos biológicos são:
(i) \emph{identificabilidade}: múltiplas combinações de parâmetros podem ajustar igualmente
bem os dados; (ii) \emph{ruído experimental}: medições são sempre imprecisas, com erros que
propagam para os parâmetros; (iii) \emph{dados escassos}: experimentos biológicos são caros
e as séries temporais são curtas; (iv) \emph{mínimos locais}: a função objetivo é
tipicamente não-convexa, com múltiplos mínimos locais que podem capturar algoritmos de
otimização; (v) \emph{correlação}: parâmetros frequentemente aparecem como produtos ou
razões no modelo, tornando impossível determiná-los individualmente.

\section{Formulação da Função Objetivo}

\subsection{Mínimos Quadrados Ordinários}

A função objetivo mais comum é a \emph{soma dos quadrados dos resíduos} (SSE, do inglês
\textit{sum of squared errors}), base do método de Mínimos Quadrados Ordinários (OLS):

\begin{equation}
J(\boldsymbol{\theta}) = \text{SSE}(\boldsymbol{\theta}) = \sum_{i=1}^{n} r_i^2 =
\sum_{i=1}^{n} \left[y_i - \hat{y}_i(\boldsymbol{\theta})\right]^2
\end{equation}

onde $r_i = y_i - \hat{y}_i(\boldsymbol{\theta})$ é o \emph{resíduo} no ponto $i$ ---
a diferença entre o dado observado e a previsão do modelo com os parâmetros atuais.
A minimização de SSE é justificada estatisticamente quando os erros de medição são
independentes e normalmente distribuídos com variância constante.

O OLS tem solução analítica para modelos lineares (a fórmula de mínimos quadrados via
álgebra linear), mas para modelos não-lineares --- como equações diferenciais ordinárias
--- a minimização requer algoritmos iterativos.

\subsection{Mínimos Quadrados Ponderados}

Quando as observações têm diferentes níveis de incerteza, o método de \emph{Mínimos
Quadrados Ponderados} (WLS) generaliza o OLS atribuindo a cada ponto um peso $w_i$
inversamente proporcional à sua variância experimental:

\begin{equation}
J(\boldsymbol{\theta}) = \sum_{i=1}^{n} w_i \left[y_i - \hat{y}_i(\boldsymbol{\theta})\right]^2,
\quad \text{tipicamente } w_i = \frac{1}{\sigma_i^2}
\end{equation}

Na cinética enzimática, por exemplo, medições a concentrações baixas de substrato tendem a
ter maior erro relativo; nesse caso, pesos $w_i = 1/\hat{y}_i$ ou $w_i = 1/\hat{y}_i^2$
são mais apropriados que pesos uniformes.

\subsection{Máxima Verossimilhança}

A abordagem de \emph{Máxima Verossimilhança} (MLE, do inglês \textit{Maximum Likelihood
Estimation}) fornece a fundamentação estatística mais rigorosa para a estimação de
parâmetros. Em vez de minimizar um erro, buscamos os parâmetros que maximizam a
probabilidade de observar os dados que de fato observamos:

\begin{equation}
\boldsymbol{\theta}^* = \arg\max_{\boldsymbol{\theta}} \mathcal{L}(\boldsymbol{\theta}) =
\arg\max_{\boldsymbol{\theta}} P(\text{dados} \mid \boldsymbol{\theta})
\end{equation}

Por conveniência numérica, trabalha-se com a \emph{log-verossimilhança}:
$\ell(\boldsymbol{\theta}) = \ln \mathcal{L}(\boldsymbol{\theta}) = \sum_i \ln P(y_i \mid \boldsymbol{\theta})$.

Quando os erros de medição são independentes e gaussianos com variância $\sigma^2$:

\begin{equation}
\ell(\boldsymbol{\theta}) = -\frac{n}{2}\ln(2\pi\sigma^2) -
\frac{1}{2\sigma^2} \sum_{i=1}^{n}\left[y_i - \hat{y}_i(\boldsymbol{\theta})\right]^2
\end{equation}

Um resultado fundamental é que, sob essa hipótese de ruído gaussiano, maximizar
$\ell(\boldsymbol{\theta})$ é matematicamente equivalente a minimizar o SSE. Ou seja:
\emph{OLS é MLE para dados com ruído gaussiano}. Quando o ruído tem outra distribuição
--- Poisson para dados de contagem, log-normal para concentrações biológicas ---
a MLE produz funções objetivo diferentes do SSE.

\subsection{Abordagem Bayesiana}

A abordagem bayesiana trata os parâmetros não como valores fixos desconhecidos, mas como
variáveis aleatórias com distribuições de probabilidade. O Teorema de Bayes conecta o
conhecimento \emph{a priori} sobre os parâmetros ($P(\boldsymbol{\theta})$) com os dados
observados para produzir uma distribuição \emph{posterior}:

\begin{equation}
P(\boldsymbol{\theta} \mid \text{dados}) = \frac{P(\text{dados} \mid \boldsymbol{\theta}) \cdot P(\boldsymbol{\theta})}{P(\text{dados})}
\end{equation}

O numerador é o produto da verossimilhança pelo prior; o denominador é uma constante de
normalização. Em vez de um único ponto ótimo, a abordagem bayesiana produz uma distribuição
completa sobre os parâmetros, quantificando incerteza de forma natural. Parâmetros bem
determinados pelos dados terão posteriors estreitos; parâmetros mal identificados terão
posteriors largos, herdando a distribuição do prior.

A principal desvantagem é computacional: calcular a posterior analiticamente é geralmente
impossível, exigindo métodos de amostragem como MCMC (\textit{Markov Chain Monte Carlo}).
Para modelos biológicos com integrações de EDOs em cada avaliação, o custo pode ser
proibitivo.


\section{Busca em Grade}

\begin{definicaoenv}[Busca em Grade]
A busca em grade (\textit{grid search}) é o método mais simples de estimação de parâmetros:
avalia sistematicamente a função objetivo em uma grade regular de valores de parâmetros,
retornando o ponto com menor $J(\boldsymbol{\theta})$.
\end{definicaoenv}

O algoritmo é direto: (1) definir limites e resolução para cada parâmetro; (2) criar uma
grade de $n_1 \times n_2 \times \cdots \times n_p$ pontos; (3) avaliar $J$ em cada nó;
(4) selecionar o ponto com menor $J$. Sua principal virtude é a garantia de encontrar o
mínimo global dentro da região explorada, desde que a grade seja suficientemente densa.

A grande limitação é o custo computacional exponencial. Com 10 valores testados por
parâmetro, 2 parâmetros requerem $10^2 = 100$ avaliações; 5 parâmetros requerem
$10^5 = 100.000$; 10 parâmetros exigiriam $10^{10}$ avaliações --- impraticável.
Esse crescimento exponencial com a dimensionalidade é chamado de \emph{maldição da
dimensionalidade} e torna a busca em grade inaplicável para mais de 2 ou 3 parâmetros.
Em prática, a busca em grade é útil para exploração inicial do espaço de parâmetros e
para gerar estimativas iniciais para algoritmos mais sofisticados.

\begin{lstlisting}[language=Python, caption=Busca em grade 2D para modelo exponencial]
import numpy as np

# Dados experimentais
t_data = np.array([0, 1, 2, 3, 4, 5])
y_data = np.array([10, 6.1, 3.7, 2.2, 1.4, 0.8])

def sse_model(params, t, y):
    a, b = params
    return np.sum((y - a * np.exp(-b * t))**2)

# Grade de busca: 20x20 pontos
a_range = np.linspace(5, 15, 20)
b_range = np.linspace(0.1, 1.0, 20)
SSE_grid = np.array([[sse_model([a, b], t_data, y_data)
                       for b in b_range] for a in a_range])

# Encontrar minimo
idx = np.unravel_index(SSE_grid.argmin(), SSE_grid.shape)
a_best, b_best = a_range[idx[0]], b_range[idx[1]]
print(f"Parametros otimos: a={a_best:.3f}, b={b_best:.3f}")
print(f"SSE minimo: {SSE_grid[idx]:.4f}")
\end{lstlisting}

\section{Métodos de Gradiente}

\subsection{Gradient Descent}

Os \emph{métodos de gradiente} exploram a geometria local da função objetivo para navegar
eficientemente em direção ao mínimo. O mais simples, \emph{gradient descent} (descida do
gradiente), move iterativamente os parâmetros na direção oposta ao gradiente --- pois o
gradiente aponta na direção de maior crescimento, e seu oposto aponta para a descida mais
íngreme:

\begin{equation}
\boldsymbol{\theta}^{(k+1)} = \boldsymbol{\theta}^{(k)} - \alpha \nabla J(\boldsymbol{\theta}^{(k)})
\end{equation}

O escalar $\alpha > 0$ é a \emph{taxa de aprendizado} (tamanho do passo). Sua escolha é
crítica: $\alpha$ muito pequeno torna a convergência lenta; $\alpha$ muito grande causa
oscilações ou divergência. Na prática, estratégias como \emph{line search} (otimizar
$\alpha$ em cada iteração) ou taxas adaptativas (Adam, RMSprop) são preferíveis ao $\alpha$
fixo.

Quando o gradiente analítico não está disponível --- o que é comum em modelos definidos por
EDOs --- usa-se aproximação por diferenças finitas:

\begin{equation}
\frac{\partial J}{\partial \theta_i} \approx \frac{J(\boldsymbol{\theta} + h \mathbf{e}_i) - J(\boldsymbol{\theta})}{h}
\end{equation}

onde $\mathbf{e}_i$ é o $i$-ésimo vetor canônico e $h \sim 10^{-5}$ é um incremento
pequeno. A convergência é declarada quando $\|\nabla J\| < \epsilon_g$, ou quando a mudança
nos parâmetros $\|\boldsymbol{\theta}^{(k+1)} - \boldsymbol{\theta}^{(k)}\| < \epsilon_\theta$,
com valores típicos $\epsilon_g = \epsilon_\theta = 10^{-6}$.

\subsection{Newton-Raphson e Gauss-Newton}

O \emph{método de Newton-Raphson} usa informação de segunda ordem --- a \emph{matriz
Hessiana} $\mathbf{H}$, com $H_{ij} = \partial^2 J / \partial \theta_i \partial \theta_j$
--- para encontrar o mínimo mais rapidamente:

\begin{equation}
\boldsymbol{\theta}^{(k+1)} = \boldsymbol{\theta}^{(k)} - \mathbf{H}^{-1}(\boldsymbol{\theta}^{(k)}) \nabla J(\boldsymbol{\theta}^{(k)})
\end{equation}

A convergência é quadrática perto do mínimo --- o número de dígitos corretos dobra a cada
iteração --- mas calcular e inverter $\mathbf{H}$ é custoso ($O(p^3)$ para $p$ parâmetros)
e o método pode divergir longe do mínimo.

Para problemas de mínimos quadrados, o \emph{método de Gauss-Newton} oferece um compromisso
elegante: aproxima a Hessiana pelo produto do \emph{Jacobiano} $\mathbf{J}$ dos resíduos
consigo mesmo ($\mathbf{H} \approx \mathbf{J}^T \mathbf{J}$, onde $J_{ij} = \partial r_i / \partial \theta_j$),
evitando o cálculo das derivadas de segunda ordem:

\begin{equation}
\boldsymbol{\theta}^{(k+1)} = \boldsymbol{\theta}^{(k)} + (\mathbf{J}^T \mathbf{J})^{-1} \mathbf{J}^T \mathbf{r}
\end{equation}

\subsection{Levenberg-Marquardt}

O algoritmo de \emph{Levenberg-Marquardt} (LM) é o método padrão para ajuste não-linear
de curvas e é a base da função \texttt{scipy.optimize.curve\_fit}. Ele interpola
adaptativamente entre Gauss-Newton (passos grandes, rápidos, perto do mínimo) e gradient
descent (passos pequenos, robustos, longe do mínimo) por meio de um parâmetro de
amortecimento $\lambda$:

\begin{equation}
\boldsymbol{\theta}^{(k+1)} = \boldsymbol{\theta}^{(k)} + (\mathbf{J}^T \mathbf{J} + \lambda \mathbf{I})^{-1} \mathbf{J}^T \mathbf{r}
\end{equation}

Quando $\lambda \to 0$, recupera-se Gauss-Newton; quando $\lambda \to \infty$, o passo se
torna um pequeno gradiente descendente. A estratégia adaptativa é simples: se um passo
reduz $J$, aceitar e diminuir $\lambda$ (tornar mais agressivo); se $J$ aumenta, rejeitar e
aumentar $\lambda$ (tornar mais conservador).

\begin{lstlisting}[language=Python, caption=Ajuste de Michaelis-Menten via Levenberg-Marquardt]
from scipy.optimize import curve_fit, least_squares
import numpy as np

def michaelis_menten(S, Vmax, Km):
    return Vmax * S / (Km + S)

# Dados experimentais de cinetica enzimatica
S_data = np.array([0.1, 0.5, 1, 2, 5, 10, 20])   # mM
v_data = np.array([0.3, 1.1, 1.8, 2.5, 3.3, 3.7, 3.9])  # umol/min

# Ajuste via LM (curve_fit usa LM internamente)
params, cov = curve_fit(michaelis_menten, S_data, v_data,
                        p0=[4.0, 2.0])
Vmax_fit, Km_fit = params
std_errors = np.sqrt(np.diag(cov))

print(f"Vmax = {Vmax_fit:.3f} +/- {std_errors[0]:.3f}")
print(f"Km   = {Km_fit:.3f}  +/- {std_errors[1]:.3f}")

# Residuos
y_pred = michaelis_menten(S_data, *params)
residuals = v_data - y_pred
print(f"RMSE = {np.sqrt(np.mean(residuals**2)):.4f}")
\end{lstlisting}


\section{Algoritmos de Otimização Global}

\subsection{O Problema dos Mínimos Locais}

Os métodos de gradiente são eficientes quando a função objetivo é convexa ou quando se
dispõe de uma boa estimativa inicial próxima do mínimo global. Em modelos biológicos reais,
porém, $J(\boldsymbol{\theta})$ é frequentemente não-convexa, com múltiplos mínimos locais
separados por barreiras altas. Um método local iniciado longe do ótimo global pode convergir
para um mínimo local, retornando parâmetros biologicamente sem sentido.

Algoritmos de otimização global contornam esse problema explorando mais amplamente o espaço
de parâmetros, com o custo de avaliações adicionais da função objetivo. Uma estratégia
híbrida --- DE (exploração global) seguido de LM (refinamento local) --- frequentemente
combina o melhor de ambos os mundos.

\subsection{Simulated Annealing}

O \emph{simulated annealing} (SA) é inspirado no processo metalúrgico de recozimento:
aquecer um metal a alta temperatura e resfriá-lo lentamente permite que os átomos escapem
de configurações de alta energia para atingir o cristal de mínima energia. No contexto de
otimização, a \emph{temperatura} $T$ controla a probabilidade de aceitar movimentos para
soluções piores:

\begin{equation}
P(\text{aceitar}) = \begin{cases}
1 & \text{se } \Delta J < 0 \\
e^{-\Delta J / T} & \text{se } \Delta J \geq 0
\end{cases}
\end{equation}

onde $\Delta J = J(\boldsymbol{\theta}_\text{novo}) - J(\boldsymbol{\theta}_\text{atual})$.
A temperatura decresce segundo um cronograma $T_k = T_0 \alpha^k$, com $\alpha \in [0.8, 0.99]$.
Em alta temperatura, a busca é amplamente exploratória; em baixa temperatura, converge para
o mínimo local mais próximo.

\subsection{Evolução Diferencial}

A \emph{Evolução Diferencial} (DE) é atualmente o algoritmo global mais recomendado para
estimação de parâmetros em biologia de sistemas. Para cada indivíduo $\boldsymbol{\theta}$
da população, um candidato mutante $\mathbf{v}$ é criado combinando três indivíduos
aleatórios $\mathbf{a}, \mathbf{b}, \mathbf{c}$:

\begin{equation}
\mathbf{v} = \mathbf{a} + F(\mathbf{b} - \mathbf{c})
\end{equation}

onde $F \in [0.5, 1.0]$ é o fator de escala de mutação. O candidato $\mathbf{u}$ é gerado
por cruzamento com probabilidade $CR \in [0.5, 0.9]$, e substitui $\boldsymbol{\theta}$
apenas se $J(\mathbf{u}) \leq J(\boldsymbol{\theta})$.

\begin{lstlisting}[language=Python, caption=Otimizacao global com Differential Evolution]
from scipy.optimize import differential_evolution
import numpy as np

# Dados experimentais
t_data = np.array([0, 1, 2, 3, 4, 5])
y_data = np.array([10, 6.1, 3.7, 2.2, 1.4, 0.8])

def objective(params, t, y):
    a, b = params
    return np.sum((y - a * np.exp(-b * t))**2)

# Limites para os parametros
bounds = [(1, 20), (0.01, 2.0)]

result = differential_evolution(
    objective, bounds, args=(t_data, y_data),
    strategy='best1bin', maxiter=1000, popsize=15,
    tol=1e-7, mutation=(0.5, 1), recombination=0.7, seed=42
)
print(f"Parametros: a={result.x[0]:.3f}, b={result.x[1]:.3f}")
print(f"SSE minimo: {result.fun:.6f}")

# Refinamento local: abordagem hibrida DE + LM
from scipy.optimize import least_squares
res_lm = least_squares(
    lambda p: y_data - p[0]*np.exp(-p[1]*t_data),
    x0=result.x, method='lm'
)
print(f"Pos-refinamento: a={res_lm.x[0]:.4f}, b={res_lm.x[1]:.4f}")
\end{lstlisting}

\subsection{Particle Swarm Optimization}

O \emph{Particle Swarm Optimization} (PSO) é inspirado no comportamento coletivo de
cardumes de peixes. Cada partícula mantém uma posição $\mathbf{x}_i$ (parâmetros) e
velocidade $\mathbf{v}_i$ no espaço de busca. A velocidade em cada iteração combina
inércia, atração pela melhor posição histórica própria $\mathbf{p}_i$ e atração pela melhor
posição global $\mathbf{g}$:

\begin{align}
\mathbf{v}_i^{(k+1)} &= w \mathbf{v}_i^{(k)} + c_1 r_1 (\mathbf{p}_i - \mathbf{x}_i^{(k)}) + c_2 r_2 (\mathbf{g} - \mathbf{x}_i^{(k)}) \\
\mathbf{x}_i^{(k+1)} &= \mathbf{x}_i^{(k)} + \mathbf{v}_i^{(k+1)}
\end{align}

onde $w$ é o peso de inércia e $r_1$, $r_2$ são números aleatórios em $[0,1]$.

A tabela a seguir compara os principais métodos de otimização para estimação de parâmetros:

\begin{table}[h]
\centering
\caption{Comparacao de metodos de otimizacao}
\label{tab:otimizacao}
\begin{tabular}{lllll}
\toprule
\textbf{Metodo} & \textbf{Escopo} & \textbf{Velocidade} & \textbf{Derivadas} & \textbf{Uso} \\
\midrule
Busca em grade & Global & Lenta & Nao & Exploracao inicial \\
Levenberg-Marquardt & Local & Rapida & Jacobiano & Padrao nao-linear \\
Simulated Annealing & Global & Lenta & Nao & Paisagens rugosas \\
Evolucao Diferencial & Global & Moderada & Nao & Parametros continuos \\
PSO & Global & Moderada & Nao & Convergencia rapida \\
\bottomrule
\end{tabular}
\end{table}


\section{Análise de Identificabilidade}

\subsection{Identificabilidade Estrutural}

Antes de iniciar qualquer estimação, é fundamental perguntar: mesmo com dados perfeitos ---
infinitos, sem ruído, medindo todas as variáveis de estado --- seria possível determinar
os parâmetros do modelo de forma única? Essa questão define a \emph{identificabilidade
estrutural}, uma propriedade do modelo em si, independente dos dados.

\begin{definicaoenv}[Identificabilidade Estrutural]
Um modelo é \emph{estruturalmente identificável} se, para dados perfeitos, existe uma
correspondência bijetiva entre parâmetros e observações --- ou seja, dois conjuntos de
parâmetros distintos produzem saídas distintas para qualquer condição inicial.
\end{definicaoenv}

Um exemplo elucidativo: o modelo $\dot{x} = k_1 - k_2 x$ tem dois parâmetros independentes,
$k_1$ e $k_2$, que podem ser determinados univocamente a partir de $x(t)$ --- é
estruturalmente identificável. Em contraste, o modelo $\dot{x} = (k_1 k_2) - k_2 x$
pode ser reescrito como $\dot{x} = k_3 - k_2 x$ com $k_3 = k_1 k_2$: apenas o produto
$k_3$ é identificável, não $k_1$ e $k_2$ separadamente. Esse modelo é \emph{estruturalmente
não-identificável}: infinitos pares $(k_1, k_2)$ com o mesmo produto $k_1 k_2$ produzirão
exatamente a mesma trajetória.

A análise de identificabilidade estrutural pode ser feita algebricamente por eliminação de
variáveis (baseada em bases de Gröbner) ou por ferramentas computacionais como DAISY,
GenSSI e STRIKE-GOLDD.

\subsection{Identificabilidade Prática}

Mesmo que um modelo seja estruturalmente identificável, pode ser praticamente impossível
determinar os parâmetros com precisão razoável a partir de dados reais, finitos e ruidosos.
Isso define a \emph{identificabilidade prática}.

As causas mais comuns de não-identificabilidade prática são: (i) dados insuficientes em
quantidade ou qualidade para discriminar entre configurações de parâmetros; (ii) parâmetros
\emph{correlacionados}, que aparecem sempre juntos nas previsões do modelo --- por exemplo,
em $y \approx a(1 - bt)$ para $t$ pequeno, os parâmetros $a$ e $ab$ são correlacionados;
(iii) baixa \emph{sensibilidade} de certas variáveis observadas a determinados parâmetros;
e (iv) \emph{superparametrização} --- mais parâmetros do que necessário para descrever
os dados disponíveis.

\subsection{Matriz de Informação de Fisher}

A \emph{Matriz de Informação de Fisher} (FIM) quantifica quanta informação os dados
contêm sobre cada parâmetro e sobre as correlações entre pares de parâmetros:

\begin{equation}
\mathbf{F}_{ij} = \sum_{k=1}^{n} \frac{1}{\sigma_k^2} \frac{\partial \hat{y}_k}{\partial \theta_i} \frac{\partial \hat{y}_k}{\partial \theta_j}
\end{equation}

A inversa da FIM fornece o \emph{limite de Cramér-Rao}: a variância de qualquer estimador
não-viesado dos parâmetros satisfaz $\text{Var}(\hat{\theta}_i) \geq (\mathbf{F}^{-1})_{ii}$.
Portanto, $\mathbf{C} = \mathbf{F}^{-1}$ é uma estimativa da \emph{matriz de covariância}
dos parâmetros. Um determinante $\det(\mathbf{F}) \approx 0$ indica não-identificabilidade
--- a FIM é (quase) singular, e a covariância é (quase) infinita.

\subsection{Intervalos de Confiança e Profile Likelihood}

Para grandes amostras, os estimadores de mínimos quadrados seguem distribuição assintótica
normal: $\hat{\boldsymbol{\theta}} \sim \mathcal{N}(\boldsymbol{\theta}^*, \mathbf{C})$.
O intervalo de confiança de 95\% para o parâmetro $\theta_i$ é:

\begin{equation}
\theta_i^* \pm t_{n-p, 0.025} \sqrt{C_{ii}}
\end{equation}

onde $t_{n-p, 0.025}$ é o quantil da distribuição $t$ de Student com $n-p$ graus de
liberdade. A diagonal $\sqrt{C_{ii}}$ fornece o erro padrão do parâmetro $\theta_i$, e os
elementos fora da diagonal indicam correlações: $\rho_{ij} = C_{ij}/\sqrt{C_{ii} C_{jj}}$
próximo de $\pm 1$ indica parâmetros altamente correlacionados.

O \emph{profile likelihood} é um método mais robusto para calcular intervalos de confiança,
especialmente quando os erros são não gaussianos ou a amostra é pequena. Para cada valor
fixo $\theta_i = \theta_i^\dagger$, otimiza-se $J$ em relação a todos os outros parâmetros
e registra-se o valor mínimo $J_\text{prof}(\theta_i^\dagger)$. O intervalo de confiança
de 95\% é o conjunto de valores de $\theta_i$ para os quais
$J_\text{prof}(\theta_i) \leq J_\text{min} + \chi^2_{1,0.95}/2 \approx J_\text{min} + 1.92$.
Um profile com forma de parábola simétrica indica parâmetro bem identificável; um profile
assintótico ou com dois mínimos indica não-identificabilidade.

\begin{lstlisting}[language=Python, caption=Calculo de intervalos de confianca e correlacao]
from scipy.optimize import curve_fit
from scipy import stats
import numpy as np

S_data = np.array([0.1, 0.5, 1, 2, 5, 10, 20])
v_data = np.array([0.3, 1.1, 1.8, 2.5, 3.3, 3.7, 3.9])

def mm(S, Vmax, Km):
    return Vmax * S / (Km + S)

params, cov = curve_fit(mm, S_data, v_data, p0=[4.0, 2.0])
std_errors = np.sqrt(np.diag(cov))

# Intervalo de confianca 95%
n, p = len(S_data), len(params)
t_val = stats.t.ppf(0.975, df=n-p)
ci = t_val * std_errors

print("Parametro | Estimativa | IC 95%")
names = ['Vmax', 'Km']
for name, val, err in zip(names, params, ci):
    print(f"{name:8s} = {val:.3f} +/- {err:.3f}")

# Matriz de correlacao
corr = cov / np.outer(std_errors, std_errors)
print(f"\nCorrelacao Vmax-Km: {corr[0,1]:.3f}")
if abs(corr[0,1]) > 0.9:
    print("ALERTA: parametros altamente correlacionados!")
\end{lstlisting}


\section{Validação e Critérios de Seleção de Modelos}

\subsection{Análise de Resíduos}

Um bom ajuste não se limita a um valor pequeno de $J$: o padrão dos resíduos
$r_i = y_i - \hat{y}_i$ é tão informativo quanto o próprio valor da função objetivo.
Quando o modelo é adequado e os erros são realmente gaussianos, os resíduos devem ser:
(i) distribuídos aleatoriamente em torno de zero, sem padrão sistemático; (ii) com
variância aproximadamente constante (homocedasticidade); e (iii) aproximadamente normais,
o que pode ser verificado por um gráfico quantil-quantil (Q-Q plot).

Desvios desse padrão revelam problemas específicos. Resíduos com tendência curva indicam
que o modelo não captura a forma funcional dos dados --- o modelo está \emph{mal
especificado}. Resíduos com variância crescente com $\hat{y}$ (padrão de funil) indicam
\emph{heterocedasticidade}, sugerindo que erros log-normais seriam mais adequados.
Resíduos com autocorrelação temporal indicam que o modelo não captura a dinâmica
completa do sistema.

\subsection{Coeficiente de Determinação}

O \emph{coeficiente de determinação} $R^2$ mede a fração da variância total dos dados
que é explicada pelo modelo:

\begin{equation}
R^2 = 1 - \frac{\text{SS}_\text{res}}{\text{SS}_\text{tot}} = 1 - \frac{\sum_i (y_i - \hat{y}_i)^2}{\sum_i (y_i - \bar{y})^2}
\end{equation}

Um $R^2$ próximo de 1 indica bom ajuste; próximo de 0 indica que o modelo não explica
os dados melhor que a média. Porém, $R^2$ tem uma limitação séria: sempre aumenta ao
adicionar parâmetros ao modelo, mesmo que esses parâmetros sejam irrelevantes. O
$R^2$ \emph{ajustado} corrige esse viés penalizando a complexidade:

\begin{equation}
R^2_\text{adj} = 1 - \frac{(1-R^2)(n-1)}{n-p-1}
\end{equation}

onde $n$ é o número de observações e $p$ o número de parâmetros. O $R^2_\text{adj}$
pode diminuir quando se adiciona parâmetros irrelevantes, tornando-o mais apropriado
para comparação entre modelos.

\subsection{Critérios de Informação: AIC e BIC}

Quando há múltiplos modelos concorrentes com diferentes complexidades, é necessário
um critério que equilibre qualidade de ajuste com parcimônia. O \emph{Critério de
Informação de Akaike} (AIC) faz essa mediação baseado na teoria da informação:

\begin{equation}
\text{AIC} = 2p - 2\ln(\mathcal{L}) \approx n\ln\!\left(\frac{\text{SSE}}{n}\right) + 2p
\end{equation}

onde $p$ é o número de parâmetros e $\mathcal{L}$ é a máxima verossimilhança. O modelo
com menor AIC é preferido. A penalidade $2p$ desencoraja modelos supercomplexos. Para
amostras pequenas ($n/p < 40$), usa-se a versão corrigida AICc:

\begin{equation}
\text{AICc} = \text{AIC} + \frac{2p(p+1)}{n-p-1}
\end{equation}

O \emph{Critério de Informação Bayesiana} (BIC) penaliza a complexidade mais
agressivamente, com $p\ln(n)$ em vez de $2p$:

\begin{equation}
\text{BIC} = p\ln(n) - 2\ln(\mathcal{L}) \approx n\ln\!\left(\frac{\text{SSE}}{n}\right) + p\ln(n)
\end{equation}

Para $n > 7$, o BIC favorece modelos mais simples que o AIC. A escolha entre AIC e BIC
reflete objetivos diferentes: o AIC otimiza a capacidade preditiva em dados novos; o BIC
tenta identificar o ``modelo verdadeiro'' subjacente. Na biologia de sistemas, onde dados
são escassos e o ruído é alto, é boa prática reportar ambos os critérios. Uma diferença
$\Delta\text{AIC} > 10$ entre dois modelos constitui evidência forte em favor do modelo
com menor AIC.

\subsection{Overfitting e Validação Cruzada}

\emph{Overfitting} ocorre quando o modelo se ajusta ao ruído dos dados de treinamento ao
invés de capturar o sinal real, resultando em excelente desempenho nos dados conhecidos mas
péssimo desempenho em dados novos. Os sintomas clássicos são: $R^2$ de treinamento muito
acima do $R^2$ de teste; intervalos de confiança excessivamente largos; e previsões fora
do domínio dos dados que divergem rapidamente. O tradeoff viés-variância é a perspectiva
teórica: modelos simples (poucos parâmetros) têm alto viés e baixa variância; modelos
complexos têm baixo viés mas alta variância.

A \emph{validação cruzada} $k$-fold (CV) avalia diretamente a capacidade preditiva do
modelo sem dados não vistos. Os dados são divididos em $k$ subconjuntos (folds); para cada
fold, o modelo é ajustado nos $k-1$ folds restantes e avaliado no fold reservado. O erro
de CV é a média dos erros em todos os folds. Valores típicos são $k=5$ ou $k=10$; quando
$k=n$, obtém-se a validação Leave-One-Out (LOO), mais custosa mas menos enviesada.

\begin{lstlisting}[language=Python, caption=Metricas de ajuste e selecao de modelos]
import numpy as np
from scipy.optimize import curve_fit

def calculate_metrics(y_true, y_pred, n_params):
    n = len(y_true)
    residuals = y_true - y_pred
    sse = np.sum(residuals**2)
    ss_tot = np.sum((y_true - np.mean(y_true))**2)

    r2 = 1 - sse/ss_tot
    r2_adj = 1 - (1-r2)*(n-1)/(n-n_params-1)
    rmse = np.sqrt(sse/n)
    aic = n*np.log(sse/n) + 2*n_params
    bic = n*np.log(sse/n) + n_params*np.log(n)
    aicc = aic + 2*n_params*(n_params+1)/(n-n_params-1)

    return {'R2': r2, 'R2_adj': r2_adj, 'RMSE': rmse,
            'AIC': aic, 'AICc': aicc, 'BIC': bic,
            'residuals': residuals}

# Comparar dois modelos: Michaelis-Menten vs Hill
S_data = np.array([0.1, 0.5, 1, 2, 5, 10, 20])
v_data = np.array([0.3, 1.1, 1.8, 2.5, 3.3, 3.7, 3.9])

def mm(S, Vmax, Km):      return Vmax*S/(Km+S)
def hill(S, Vmax, Km, n): return Vmax*S**n/(Km**n+S**n)

p1, _ = curve_fit(mm,   S_data, v_data, p0=[4, 2])
p2, _ = curve_fit(hill, S_data, v_data, p0=[4, 2, 1.5])

m1 = calculate_metrics(v_data, mm(S_data, *p1), 2)
m2 = calculate_metrics(v_data, hill(S_data, *p2), 3)

print("Modelo MM (2 param): AICc={:.2f}, BIC={:.2f}".format(m1['AICc'], m1['BIC']))
print("Modelo Hill (3 par): AICc={:.2f}, BIC={:.2f}".format(m2['AICc'], m2['BIC']))
print("dAICc = {:.2f}".format(m2['AICc'] - m1['AICc']))
\end{lstlisting}


\section{Exercícios}

\begin{exercicio}[Busca em Grade]
Implemente busca em grade para ajustar o modelo de potência $y = a x^b$ aos dados abaixo.
Use uma grade $20 \times 20$ com $a \in [1, 3]$ e $b \in [1.5, 2.5]$. Plote o mapa de
contorno de SSE no espaço $(a, b)$ e identifique visualmente o mínimo. Compare o resultado
com o obtido por \texttt{scipy.optimize.curve\_fit}.

\begin{center}
\begin{tabular}{c|cccccc}
$x$ & 1 & 2 & 3 & 4 & 5 & 6 \\
\hline
$y$ & 2.1 & 8.3 & 18.2 & 32.5 & 50.1 & 72.3
\end{tabular}
\end{center}
\end{exercicio}

\begin{exercicio}[Gradient Descent]
Implemente gradient descent manualmente para minimizar $J(\theta) = (\theta - 5)^2 + 3$,
usando $\theta_0 = 0$ como ponto inicial. Teste quatro valores de taxa de aprendizado:
$\alpha \in \{0.01, 0.1, 0.5, 1.0\}$. Para cada $\alpha$, plote a trajetória de $J$ ao
longo das iterações e registre o número de iterações até convergência. Qual é o valor ideal?
Explique o que acontece com $\alpha = 1.0$.
\end{exercicio}

\begin{exercicio}[Levenberg-Marquardt]
Ajuste os parâmetros da equação de Michaelis-Menten $v = V_\text{max}[S]/(K_M + [S])$
aos dados abaixo usando \texttt{scipy.optimize.curve\_fit} (algoritmo LM). Calcule os
intervalos de confiança de 95\%, a matriz de correlação entre $V_\text{max}$ e $K_M$, e
plote: (a) os dados com a curva ajustada; (b) o gráfico de resíduos vs. valores preditos.

\begin{center}
\begin{tabular}{c|ccccccc}
$[S]$ (mM) & 0.1 & 0.5 & 1 & 2 & 5 & 10 & 20 \\
\hline
$v$ ($\mu$mol/min) & 0.3 & 1.1 & 1.8 & 2.5 & 3.3 & 3.7 & 3.9
\end{tabular}
\end{center}
\end{exercicio}

\begin{exercicio}[Differential Evolution]
Repita o exercício anterior usando \texttt{scipy.optimize.differential\_evolution} com
limites amplos $V_\text{max} \in [0.1, 10]$ e $K_M \in [0.1, 10]$. Em seguida, use
o resultado do DE como estimativa inicial para LM (abordagem híbrida). Compare o tempo
de execução e o SSE final dos três métodos: LM puro, DE puro e DE+LM.
\end{exercicio}

\begin{exercicio}[Crescimento Logístico]
Gere dados sintéticos de crescimento bacteriano usando o modelo logístico
$N(t) = K / \{1 + [(K/N_0) - 1] e^{-rt}\}$ com $N_0 = 1$, $r = 0.5$, $K = 100$,
$t \in [0, 20]$ em 15 pontos equidistantes, adicionando ruído gaussiano de 5\%.
(a) Estime $r$ e $K$ por três métodos diferentes. (b) Calcule o profile likelihood
para cada parâmetro e plote o perfil vs. o limiar de confiança de 95\%.
(c) Avalie identificabilidade: $r$ e $K$ são bem identificáveis nesses dados?
\end{exercicio}

\begin{exercicio}[Lotka-Volterra --- Projeto Integrador]
Considere o sistema predador-presa de Lotka-Volterra:
\begin{align*}
\frac{dV}{dt} &= \alpha V - \beta V P \\
\frac{dP}{dt} &= \delta \beta V P - \gamma P
\end{align*}
Gere dados sintéticos simulando o sistema com $\alpha = 1.0$, $\beta = 0.1$,
$\delta = 0.75$, $\gamma = 1.5$, condições iniciais $V(0) = 10$, $P(0) = 5$,
intervalo $[0, 30]$ com ruído gaussiano de 10\%.

Em seguida: (a) formule a função SSE usando ambas as variáveis $V$ e $P$; (b) estime os
quatro parâmetros por DE; (c) refine por LM; (d) calcule e interprete a matriz de
correlação de Fisher; (e) calcule intervalos de confiança de 95\% e avalie quais
parâmetros são bem identificáveis; (f) simule o modelo com os parâmetros estimados
e compare com os dados originais.
\end{exercicio}

\section{Notas Bibliográficas}

O problema de estimação de parâmetros em modelos biológicos tem extensa literatura.
Moles et al.\ (2003) \cite{moles2003} realizaram o primeiro estudo comparativo sistemático
de algoritmos de otimização global para sistemas bioquímicos, concluindo que algoritmos
evolutivos como DE são superiores aos métodos locais na maioria dos cenários. Raue et al.\
(2009) \cite{raue2009} introduziram a análise de profile likelihood no contexto de sistemas
biológicos, estabelecendo o método como padrão para análise de identificabilidade prática.

Para a teoria de otimização numérica em geral, Nocedal e Wright \cite{nocedal2006} é a
referência definitiva, cobrindo métodos de gradiente, Newton e quasi-Newton com rigor
matemático. Press et al.\ \cite{press2007} (\textit{Numerical Recipes}) apresentam
implementações práticas comentadas.

O fenômeno de ``sloppy models'' --- modelos biológicos nos quais certas combinações de
parâmetros são bem determinadas pelos dados enquanto outras são essencialmente
inidentificáveis --- foi estudado por Gutenkunst et al.\ (2007) \cite{gutenkunst2007}.
Esse trabalho sugere que a não-identificabilidade é uma característica ubíqua de modelos
de sistemas biológicos, não uma exceção.

Para seleção de modelos, Burnham e Anderson \cite{burnham2002} é a referência canônica
sobre AIC e AICc, com vasta discussão sobre inferência multimodelo. Para abordagens
bayesianas, o pacote PyMC (pymc.io) oferece uma implementação acessível de MCMC para
modelos com integração de EDOs.

\begin{thebibliography}{99}

\bibitem{moles2003}
Moles C.G., Mendes P., Banga J.R. (2003).
Parameter estimation in biochemical pathways: a comparison of global optimization methods.
\textit{Genome Research} \textbf{13}: 2467--2474.

\bibitem{raue2009}
Raue A. et al. (2009).
Structural and practical identifiability analysis of partially observed dynamical models
by exploiting the profile likelihood.
\textit{Bioinformatics} \textbf{25}: 1923--1929.

\bibitem{nocedal2006}
Nocedal J., Wright S.J. (2006).
\textit{Numerical Optimization}, 2ª ed.
Springer.

\bibitem{press2007}
Press W.H., Teukolsky S.A., Vetterling W.T., Flannery B.P. (2007).
\textit{Numerical Recipes: The Art of Scientific Computing}, 3ª ed.
Cambridge University Press.

\bibitem{gutenkunst2007}
Gutenkunst R.N. et al. (2007).
Universally sloppy parameter sensitivities in systems biology models.
\textit{PLoS Computational Biology} \textbf{3}: e189.

\bibitem{burnham2002}
Burnham K.P., Anderson D.R. (2002).
\textit{Model Selection and Multimodel Inference: A Practical Information-Theoretic Approach},
2ª ed. Springer.

\bibitem{voit2015}
Voit E.O., Martens H.A., Omholt S.W. (2015).
150 years of the mass action law.
\textit{PLoS Computational Biology} \textbf{11}: e1004012.

\end{thebibliography}
